\subsection{Bar–cobar construction}
\label{p:sec:bar-cobar}


Chiaffrino, Mamade and Zwiebach have shown that the equations of motion of any $\Loo$ field theory, closed string field theory included, can be rewritten as a quadratic equation on an enlarged space of multi-string fields~\cite{Chiaffrino:2026:UniversalQuadraticField}.
Their proofs are written for $A_\infty$ algebras. For $\Loo$ algebras, the equivalence follows from a general theorem on quasi-isomorphic algebras.
They give no action for that equation and leave its existence open.
The enlarged space contains one multi-string field for every number of strands, which they call type I, besides fields valued in formal products of such multi-string spaces, their type II (both are defined below).
Its $2$-strand component is the space in which $\Sigma_2$ lives, up to a quotient, so that we can compare the two constructions field by field.
In this section, we recall their construction in our notation, we identify the fields and the equations which the two constructions share, and we state in what sense neither construction is a truncation of the other.
The comparison is classical and algebraic.


\subsubsection{The bar–cobar construction}
\label{p:sec:barcobar:construction}

The products $\ell_n$ of closed string field theory form an $\Loo$ algebra~\cite{Zwiebach:1993:ClosedStringField,Erbin:2021:StringFieldTheory}.
The bar construction packages all of them into a single coderivation $B$ of the symmetric tensor coalgebra built on the state space, and the cobar construction then builds the free Lie algebra on this coalgebra.
The result is a differential graded Lie algebra, quasi-isomorphic to the original $\Loo$ algebra~\cite{Chiaffrino:2026:UniversalQuadraticField}.
Its Maurer--Cartan equation is quadratic because a differential graded Lie algebra has no product beyond the bracket. The price is the size of the field space.

The field $\Phi$ of the enlarged theory has components labeled by a \emph{factor number} $f$ and by a \emph{length}, the number of strands.
The components of factor number one, which they call type I, are one multi-string field per length:
\begin{equation}
  \label{p:eq:barcobar:type-I}
  \Phi
    = \Phi_1 + \Phi_2 + \cdots,
  \qquad
  \Phi_1
    = \sum_{r \ge 1} \Phi_{1, r},
  \qquad
  \Phi_{1, r}
    \in \Hn{r},
  \qquad
  \Ngh(\Phi_{1, r})
    = 2 r,
\end{equation}
and the length-one component $\Phi_{1,1}$ is the string field itself.
We describe their construction in our notation, translating their homological degree into ghost number: a type-I field of length $r$ and total ghost number $N$ has degree
\begin{equation}
  \label{p:eq:barcobar:degree}
  \deg \Phi_{1, r}
    = 2 r - 1 - N
\end{equation}
in their enlarged space.
The components of factor number $f \ge 2$, their type II, are independent fields valued in the formal products $\Hn{r_1} \otimes' \cdots \otimes' \Hn{r_f}$ generated by the free Lie algebra of the cobar step.
A type-II component with $f$ factors and total length $r$ has ghost number $2 r + 1 - f$, which differs from the ghost number $2 r$ of the auxiliary field with the same number of strands.
The equation of motion reads:
\begin{equation}
  \label{p:eq:barcobar:mc}
  (B + \Delta_c) \Phi + \Phi \otimes' \Phi
    = 0,
\end{equation}
where $\Phi \otimes' \Phi = \tfrac{1}{2} [\Phi, \Phi]$ for a field of their degree $-1$.
All the data of the theory sit in $B$: on a type-I field of length $r$ it applies the products $\gs^{n-1} \ell_n$, for every $n \ge 1$, to $n$ of the strands and sums over the choice of the strands.
It lowers the length by $n - 1$ and preserves the factor number.
The two other operations are universal and raise the factor number by one.
The coproduct $\Delta_c$ splits a type-I field of length $n$ into the pairs of lengths $(i, n - i)$, by summing over the ways of distributing the $n$ strands, and $\otimes'$ juxtaposes two fields.

Since $B$ preserves the factor number and the two other operations raise it by one, an ansatz with type-I components only, $\Phi = \Phi_1$, reduces~\eqref{p:eq:barcobar:mc} to two equations:
\begin{equation}
  \label{p:eq:barcobar:canonical-eqs}
  B \Phi_1
    = 0,
  \qquad
  \Delta_c \Phi_1 + \Phi_1 \otimes' \Phi_1
    = 0,
\end{equation}
of factor numbers one and two, all higher factor numbers vanishing identically.
The second equation is a factorization condition: component by component, for every length $n$ and every split $1 \le i \le n - 1$:
\begin{equation}
  \label{p:eq:barcobar:factorization}
  \binom{n}{i} \, \Phi_{1, n}
    = \Phi_{1, i} \otimes \Phi_{1, n - i},
\end{equation}
as tensors with $n$ strands, the right-hand side being required to be symmetric under all permutations of its strands.
Its unique solution is the exponential of the string field:
\begin{equation}
  \label{p:eq:barcobar:canonical}
  \Phi_{1, r}
    = \frac{\Psi^r}{r!},
\end{equation}
which satisfies~\eqref{p:eq:barcobar:factorization} for every split.
On this solution, the first equation of~\eqref{p:eq:barcobar:canonical-eqs} reduces to the equation of motion $\mathcal E_\Psi = 0$ of the original theory, with $\mathcal E_\Psi$ given by~\eqref{p:eq:hs1:eom-psi}.
Indeed, the components of $B \Phi_1$ with $L \ge 2$ output strands are $\mathcal E_\Psi$ tensored with $\Psi^{L-1} / (L-1)!$.

They call these solutions \emph{canonical}.
They prove, order by order in $B$, that every solution of~\eqref{p:eq:barcobar:mc} is gauge equivalent to a canonical one~\cite{Chiaffrino:2026:UniversalQuadraticField}.
The gauge classes of solutions of the quadratic equation are therefore in one-to-one correspondence with those of the original theory.
They also show that, among the parameters of factor number one, those which preserve the canonical form generate the gauge transformations of the original theory.
In the $\Loo$ case, where the bar construction uses the symmetric coalgebra, the same argument gives a single-string parameter $\Lambda \in \Hs$ of ghost number one.

Their construction gives equations but no action, and they identify two obstructions~\cite{Chiaffrino:2026:UniversalQuadraticField}.
First, an action for a cyclic $\Loo$ algebra needs a non-degenerate bilinear form of a definite degree, and they argue that the enlarged space admits none.
In their scalar field example every vector has negative degree, so that nothing is available to pair against.
Second, they argue that a form diagonal in the factor number and the length is unlikely to be cyclic, since $\Delta_c$ and $\otimes'$ raise the factor number.
Indeed $\langle \Phi_2, \Phi_1 \otimes' \Phi_1 \rangle$ can be nonzero while $\langle \Phi_1, \Phi_2 \otimes' \Phi_1 \rangle$ vanishes, and for the same reason $\Delta_c$ is unlikely to be self-adjoint.
They leave the question open.

We can easily write an action whose equations contain theirs by doubling every field with a multiplier. The question they leave open is the existence of an action on their field space itself.
They also note that simpler constructions should be possible with more knowledge of the theory, and give the Hubbard--Stratonovich transformation of a quartic interaction with one bilocal field as an example.
For closed string field theory this example is the earlier construction of~\cite{Erbin:2018:CubicClosedString}, which the paper extends to all orders.

\subsubsection{Identification on the canonical sector}
\label{p:sec:barcobar:canonical}

The type-I component of length $r$ lives in $\Hn{r}$ with ghost number $2 r$, which is the space and the ghost number of the auxiliary field $\Sigma_r$.
The identification is simplest when we absorb the factorials:
\begin{equation}
  \label{p:eq:barcobar:identification}
  \Sigma_r
    := r! \, \Phi_{1, r},
  \qquad
  r \ge 1,
\end{equation}
so that $\Sigma_1 = \Phi_{1, 1} = \Psi$, the canonical solution is $\Sigma_r = \Psi^r$, and the factorization conditions~\eqref{p:eq:barcobar:factorization} read $\Sigma_n = \Sigma_i \otimes \Sigma_{n-i}$ for every split.
